\chapter*{Abstract}
\addcontentsline{toc}{chapter}{Abstract}

The volume of artificial-intelligence-related content, spanning research papers, open-source code releases, industry announcements, government policy, funding news, and long-form video, has grown far faster than any individual can read, leaving readers unable to reliably separate what matters to them from background noise. AI Compass addresses this problem of information overload with a personalized content-intelligence platform built around a strict separation between shared, one-time computation and per-user, read-time computation. An offline ingestion pipeline continuously collects content from sixteen registered sources spanning research, open-source releases, video, developer communities, government publications, and funding news, plus user-submitted feeds screened by an automated embedding-similarity relevance check. Each item is processed once through a single structured large-language-model call that extracts a summary, controlled-vocabulary topics, named entities, a content category, and a technical-depth score, replacing what would otherwise require several separate calls per item. Content is embedded locally into 384-dimensional vectors, enabling near-duplicate story clustering, semantic search across articles and video, and passage-level retrieval. A deterministic, non-generative ranking service then scores and orders content per user from five weighted signals (interest affinity, quality, freshness, source affinity, and novelty), diversified through Maximal Marginal Relevance and a reserved exploration slice, with every recommendation paired with a plain-language, template-generated explanation. Behavioral signals such as clicks, saves, and follows are aggregated into decaying per-topic, per-source, and per-entity affinities and a personal taste vector that continuously refine this ranking. A retrieval-augmented generation (RAG) conversational assistant answers user questions strictly from indexed source passages, with every citation validated server-side before being shown, and long videos are transcribed locally and chaptered for deep summaries. The system is deployed as a split-service architecture, with the frontend hosted on a managed platform and the backend, task queues, and database self-hosted on a single cloud virtual machine. An offline evaluation harness implementing NDCG@$k$ and Mean Average Precision was implemented, methodologically corrected, and executed against the live database; it confirms the harness computes correctly but also demonstrates, with a concrete root-cause analysis, that the platform's current interaction volume is not yet sufficient to produce a statistically meaningful ranking-quality number: a limitation reported explicitly rather than concealed. Verification confirmed that the ranking pipeline produces a complete, explained, personalized feed with no large-language-model calls at all in its serving path, and that local speech-to-text transcription sustains roughly 1.76 times real-time throughput on CPU hardware alone. These results show that a personalized recommendation and conversational system for a fast-moving technical domain can confine generative models to well-defined, independently verifiable roles rather than to the ranking or retrieval core itself.

\bigskip
\noindent\textbf{Keywords:} personalized recommendation, retrieval-augmented generation, content intelligence, semantic search, large language models, information retrieval evaluation.
